\documentclass[10pt,a4paper]{article}
\usepackage[utf8]{inputenc} % para poder usar tildes en archivos UTF-8
\usepackage[spanish]{babel} % para que comandos como \today den el resultado en castellano
\usepackage{a4wide} % márgenes un poco más anchos que lo usual
\usepackage{caratula}
\usepackage{verbatimbox}
%\usepackage{algorithm}
%\usepackage{algorithmic}
%\usepackage{algpseudocode}
\usepackage{float} %para decir a donde van los pseudocodigos
\usepackage{icomma}
\usepackage{multicol}
\usepackage{todonotes}
\usepackage[nottoc]{tocbibind}
\usepackage{amsthm}
\usepackage{amssymb} %Simbolos matematicos
\usepackage{amsmath}
\usepackage{mathtools}
\usepackage{tikz,pgf}
\usetikzlibrary{babel,arrows,automata,graphs,graphs.standard,quotes,decorations.markings} %Para hacer grafos con tikz
\usepackage{multirow} %Para las tablas

\usepackage{rotating} %Para escribir vertical


%\alglanguage{pseudocode}


\newcommand{\bigO}{\mathcal{O}}

\newcommand{\grafiquito}[2]{%
    \begin{figure}[H]
    \centering
    \includegraphics[width=7cm]{./#1.png}
    \caption{#2}
    \label{fig:#1}
    \end{figure}%
}

\newcommand{\png}[2]{%
    \begin{figure}[H]
    \centering
    \includegraphics[width=\textwidth]{./img/#1.png}
    \caption{#2}
    \label{png:#1}
    \end{figure}%
}

\newcommand{\jpg}[2]{%
    \begin{figure}[H]
    \centering
    \includegraphics[width=\textwidth]{./img/#1.jpg}
    \caption{#2}
    \label{jpg:#1}
    \end{figure}%
}

\tikzset{
  big arrow/.style={
    decoration={markings,mark=at position 1 with {\arrow[scale=4,#1]{>}}},
    postaction={decorate},
    shorten >=0.4pt},
  big arrow/.default=black}
  
\newtheorem{obs}{Observaci\'on}[section]
\newtheorem{lema}{Lema}[section]
\newtheorem{ejemplo}{Ejemplo}[section]
\renewcommand*{\proofname}{Demostraci\'on}


\begin{document}

% **************************************************************************
%
%  Package 'caratula', version 0.5 (para componer caratulas de TPs del DC).
%
%  En caso de dudas, problemas o sugerencias sobre este package escribir a
%  Brian J. Cardiff (bcardif arroba gmail.com).
%  Nico Rosner (nrosner arroba dc.uba.ar).
%
% **************************************************************************

% ----- Informacion sobre el package para el sistema -----------------------

\NeedsTeXFormat{LaTeX2e}
\ProvidesPackage{caratula}[2013/08/04 v0.5 Para componer caratulas de TPs del DC]
\RequirePackage{ifthen}
\usepackage[pdftex]{graphicx}

% ----- Imprimir un mensajito al procesar un .tex que use este package -----

\typeout{Cargando package 'caratula' v0.5 (2013/08/04)}

% ----- Algunas variables --------------------------------------------------

\let\Materia\relax
\let\Submateria\relax
\let\Titulo\relax
\let\Subtitulo\relax
\let\Grupo\relax
\let\Fecha\relax
\let\Logoimagefile\relax
\newcommand{\LabelIntegrantes}{}
\newboolean{showLU}
\newboolean{showEntregas}
\newboolean{showDirectores}
\newboolean{showCoDirectores}

% ----- Comandos para que el usuario defina las variables ------------------

\def\materia#1{\def\Materia{#1}}
\def\submateria#1{\def\Submateria{#1}}
\def\titulo#1{\def\Titulo{#1}}
\def\subtitulo#1{\def\Subtitulo{#1}}
\def\grupo#1{\def\Grupo{#1}}
\def\fecha#1{\def\Fecha{#1}}
\def\logoimagefile#1{\def\Logoimagefile{#1}}

% ----- Token list para los integrantes ------------------------------------

\newtoks\intlist\intlist={}

\newtoks\intlistSinLU\intlistSinLU={}

\newcounter{integrantesCount}
\setcounter{integrantesCount}{0}
\newtoks\intTabNombre\intTabNombre={}
\newtoks\intTabLU\intTabLU={}
\newtoks\intTabEmail\intTabEmail={}

\newcounter{directoresCount}
\setcounter{directoresCount}{0}
\newtoks\direcTabNombre\direcTabNombre={}
\newtoks\direcTabEmail\direcTabEmail={}

\newcounter{coDirectoresCount}
\setcounter{coDirectoresCount}{0}
\newtoks\codirecTabNombre\codirecTabNombre={}
\newtoks\codirecTabEmail\codirecTabEmail={}


% ----- Comando para que el usuario agregue integrantes --------------------

\def\integrante#1#2#3{%
    \intlist=\expandafter{\the\intlist\rule{0pt}{1.2em}#1&#2&\tt #3\\[0.2em]}%
    \intlistSinLU=\expandafter{\the\intlistSinLU\rule{0pt}{1.2em}#1 & \tt #3\\[0.2em]}%
    %
    \ifthenelse{\value{integrantesCount} > 0}{%
        \intTabNombre=\expandafter{\the\intTabNombre & #1}%
        \intTabLU=\expandafter{\the\intTabLU & #2}%
        \intTabEmail=\expandafter{\the\intTabEmail & \tt #3}%
    }{
        \intTabNombre=\expandafter{\the\intTabNombre #1}%
        \intTabLU=\expandafter{\the\intTabLU #2}%
        \intTabEmail=\expandafter{\the\intTabEmail \tt #3}%
    }%
    \addtocounter{integrantesCount}{1}%
}

\def\director#1#2{%
    \ifthenelse{\value{directoresCount} > 0}{%
        \direcTabNombre=\expandafter{\the\direcTabNombre & #1}%
        \direcTabEmail=\expandafter{\the\direcTabEmail & \tt #2}%
    }{
        \direcTabNombre=\expandafter{\the\direcTabNombre #1}%
        \direcTabEmail=\expandafter{\the\direcTabEmail \tt #2}%
    }%
    \addtocounter{directoresCount}{1}%
}

\def\codirector#1#2{%
    \ifthenelse{\value{coDirectoresCount} > 0}{%
        \codirecTabNombre=\expandafter{\the\codirecTabNombre & #1}%
        \codirecTabEmail=\expandafter{\the\codirecTabEmail & \tt #2}%
    }{
        \codirecTabNombre=\expandafter{\the\codirecTabNombre #1}%
        \codirecTabEmail=\expandafter{\the\codirecTabEmail \tt #2}%
    }%
    \addtocounter{coDirectoresCount}{1}%
}


% ----- Macro para generar la tabla de integrantes -------------------------

\newcommand{\tablaIntegrantes}{\ }

\newcommand{\tablaIntegrantesVertical}{%
\ifthenelse{\boolean{showLU}}{%
    \begin{tabular}[t]{| l @{\hspace{4ex}} c @{\hspace{4ex}} l|}
        \hline
        \multicolumn{1}{|c}{\rule{0pt}{1.2em} \LabelIntegrantes} & LU &  \multicolumn{1}{c|}{Correo electr\'onico} \\[0.2em]
        \hline \hline
        \the\intlist
        \hline
    \end{tabular}
}{
    \begin{tabular}[t]{| l @{\hspace{4ex}} @{\hspace{4ex}} l|}
        \hline
        \multicolumn{1}{|c}{\rule{0pt}{1.2em} \LabelIntegrantes} &  \multicolumn{1}{c|}{Correo electr\'onico} \\[0.2em]
        \hline \hline
        \the\intlistSinLU
        \hline
    \end{tabular}
    }%
}

\newcommand{\tablaIntegrantesHorizontal}{%
    \begin{tabular}[t]{ *{\value{integrantesCount}}{c} }
    \the\intTabNombre \\%
\ifthenelse{\boolean{showLU}}{
    \the\intTabLU \\%
}{}
    \the\intTabEmail %
    \end{tabular}%
}

\newcommand{\tablaDirectores}{%
\ifthenelse{\boolean{showDirectores}}{%
	\bigskip
	Directores

	\smallskip
    \begin{tabular}[t]{ *{\value{directoresCount}}{c} }
    \the\direcTabNombre \\%
    \the\direcTabEmail %
    \end{tabular}%
}{}%
}

\newcommand{\tablaCoDirectores}{%
\ifthenelse{\boolean{showCoDirectores}}{%
	\bigskip
	Co-Directores

	\smallskip
    \begin{tabular}[t]{ *{\value{coDirectoresCount}}{c} }
    \the\codirecTabNombre \\%
    \the\codirecTabEmail %
    \end{tabular}%
}{}%
}

\newcommand{\tablaEntregas}{%
\ifthenelse{\boolean{showEntregas}}{%
  \bigskip%
  \begin{tabular}[t]{|l p{3.5cm} p{1.5cm}|}%
  \hline%
  \rule{0pt}{1.2em} Instancia & Docente & Nota \\[0.2em] %
  \hline%
  \hline%
  \rule{0pt}{1.2em} Primera entrega & & \\[0.2em] %
  \hline%
  \rule{0pt}{1.2em} Segunda entrega & & \\[0.2em] %
  \hline%
  \end{tabular}%
}{}%
}

% ----- Codigo para manejo de errores --------------------------------------

\def\se{\let\ifsetuperror\iftrue}
\def\ifsetuperror{%
    \let\ifsetuperror\iffalse
    \ifx\Materia\relax\se\errhelp={Te olvidaste de proveer una \materia{}.}\fi
    \ifx\Titulo\relax\se\errhelp={Te olvidaste de proveer un \titulo{}.}\fi
    \edef\mlist{\the\intlist}\ifx\mlist\empty\se%
    \errhelp={Tenes que proveer al menos un \integrante{nombre}{lu}{email}.}\fi
    \expandafter\ifsetuperror}

\def\aftermaketitle{%
  \setcounter{page}{1}
}

% ----- \maketitletxt correspondiente a la versi�n v0.2.1 (texto v0.2 + fecha ) ---------

\def\maketitletxt{%
    \ifsetuperror\errmessage{Faltan datos de la caratula! Ingresar 'h' para mas informacion.}\fi
    \thispagestyle{empty}
    \begin{center}
    \vspace*{\stretch{2}}
    {\LARGE\textbf{\Materia}}\\[1em]
    \ifx\Submateria\relax\else{\Large \Submateria}\\[0.5em]\fi
    \ifx\Fecha\relax\else{\Large \Fecha}\\[0.5em]\fi
    \par\vspace{\stretch{1}}
    {\large Departamento de Computaci\'on}\\[0.5em]
    {\large Facultad de Ciencias Exactas y Naturales}\\[0.5em]
    {\large Universidad de Buenos Aires}
    \par\vspace{\stretch{3}}
    {\Large \textbf{\Titulo}}\\[0.8em]
    {\Large \Subtitulo}
    \par\vspace{\stretch{3}}
    \ifx\Grupo\relax\else\textbf{\Grupo}\par\bigskip\fi
    \tablaIntegrantes
    \end{center}
    \vspace*{\stretch{3}}
    \newpage\aftermaketitle}

% ----- \maketitletxtlogo correspondiente v0.2.1 (texto con fecha y logo) ---------

\def\maketitletxtlogo{%
    \ifsetuperror\errmessage{Faltan datos de la caratula! Ingresar 'h' para mas informacion.}\fi
    \thispagestyle{empty}
    \begin{center}
    \ifx\Logoimagefile\relax\else\includegraphics{\Logoimagefile}\fi \hfill \includegraphics{logo_dc.jpg}\\[1em]
    \vspace*{\stretch{2}}
    {\LARGE\textbf{\Materia}}\\[1em]
    \ifx\Submateria\relax\else{\Large \Submateria}\\[0.5em]\fi
    \ifx\Fecha\relax\else{\large \Fecha}\\[0.5em]\fi
    \par\vspace{\stretch{1}}
    {\large Departamento de Computaci\'on}\\[0.5em]
    {\large Facultad de Ciencias Exactas y Naturales}\\[0.5em]
    {\large Universidad de Buenos Aires}
    \par\vspace{\stretch{3}}
    {\Large \textbf{\Titulo}}\\[0.8em]
    {\Large \Subtitulo}
    \par\vspace{\stretch{3}}
    \ifx\Grupo\relax\else\textbf{\Grupo}\par\bigskip\fi
    \tablaIntegrantes
    \end{center}
    \vspace*{\stretch{4}}
    \newpage\aftermaketitle}

% ----- \maketitlegraf correspondiente a la versi�n v0.3 (gr�fica) -------------

\def\maketitlegraf{%
    \ifsetuperror\errmessage{Faltan datos de la caratula! Ingresar 'h' para mas informacion.}\fi
%
    \thispagestyle{empty}

    \ifx\Logoimagefile\relax\else\includegraphics{\Logoimagefile}\fi \hfill \includegraphics{logo_dc.jpg}

    \vspace*{.06 \textheight}

    \noindent \textbf{\huge \Titulo}  \medskip \\
    \ifx\Subtitulo\relax\else\noindent\textbf{\large \Subtitulo} \\ \fi%
    \noindent \rule{\textwidth}{1 pt}

    {\noindent\large\Fecha \hspace*\fill \Materia} \\
    \ifx\Submateria\relax\else{\noindent \hspace*\fill \Submateria}\fi%

    \medskip%
    \begin{center}
        \ifx\Grupo\relax\else\textbf{\Grupo}\par\bigskip\fi
        \tablaIntegrantes

        \tablaDirectores

        \tablaCoDirectores

        \tablaEntregas
    \end{center}%
    \vfill%
%
    \begin{minipage}[t]{\textwidth}
        \begin{minipage}[t]{.55 \textwidth}
            \includegraphics{logo_uba.jpg}
        \end{minipage}%%
        \begin{minipage}[b]{.45 \textwidth}
            \textbf{\textsf{Facultad de Ciencias Exactas y Naturales}} \\
            \textsf{Universidad de Buenos Aires} \\
            {\scriptsize %
            Ciudad Universitaria - (Pabell\'on I/Planta Baja) \\
                Intendente G\"uiraldes 2610 - C1428EGA \\
            Ciudad Aut\'onoma de Buenos Aires - Rep. Argentina \\
                Tel/Fax: (++54 +11) 4576-3300 \\
            http://www.exactas.uba.ar \\
            }
        \end{minipage}
    \end{minipage}%
%
    \newpage\aftermaketitle}

% ----- Reemplazamos el comando \maketitle de LaTeX con el nuestro ---------
\renewcommand{\maketitle}{\maketitlegraf}

% ----- Dependiendo de las opciones ---------
%
% opciones:
%   txt     : caratula solo texto.
%   txtlogo : caratula txt con logo del DC y del grupo (opcional).
%   graf    : (default) caratula grafica con logo del DC, UBA y del grupo (opcional).
%
\@makeother\*% some package redefined it as a letter (as color.sty)
%
% Layout general de la caratula
%
\DeclareOption{txt}{\renewcommand{\maketitle}{\maketitletxt}}
\DeclareOption{txtlogo}{\renewcommand{\maketitle}{\maketitletxtlogo}}
\DeclareOption{graf}{\renewcommand{\maketitle}{\maketitlegraf}}
%
% Etiqueta Autores o Integrantes
%
\DeclareOption{integrante}{\renewcommand{\LabelIntegrantes}{Integrante}}
\DeclareOption{autor}{\renewcommand{\LabelIntegrantes}{Autor}}
%
% Formato tabla de integrantes
%
\DeclareOption{intVert}{\renewcommand{\tablaIntegrantes}{\tablaIntegrantesVertical}}
\DeclareOption{intHoriz}{\renewcommand{\tablaIntegrantes}{\tablaIntegrantesHorizontal}}
\DeclareOption{conLU}{\setboolean{showLU}{true}}
\DeclareOption{sinLU}{\setboolean{showLU}{false}}
\DeclareOption{conEntregas}{\setboolean{showEntregas}{true}}
\DeclareOption{sinEntregas}{\setboolean{showEntregas}{false}}
\DeclareOption{showDirectores}{\setboolean{showDirectores}{true}}
\DeclareOption{hideDirectores}{\setboolean{showDirectores}{false}}
\DeclareOption{showCoDirectores}{\setboolean{showCoDirectores}{true}}
\DeclareOption{hideCoDirectores}{\setboolean{showCoDirectores}{false}}
%
% Opciones predeterminadas
%
\ExecuteOptions{intVert}%
\ExecuteOptions{graf}%
\ExecuteOptions{integrante}%
\ExecuteOptions{conLU}%
\ExecuteOptions{hideDirectores}%
\ExecuteOptions{hideCoDirectores}%
\ExecuteOptions{sinEntregas}%
%
\ProcessOptions\relax


\newpage
\begin{center}
\end{center}
\thispagestyle{empty}
\setcounter{page}{0}

\newpage
\tableofcontents

% \twocolumn

\newpage
\section{Introducción}

En la actualidad existen diversas técnicas para determinar la tabla de posiciones en una competencia compuesta por varios equipos. En este trabajo vamos a estudiar varias de ellas, principalmente el método de la Mátriz de Colley (CMM), elaborado por Wes Colley.

Primero, debemos entender cuál es el problema a resolver.

\noindent \textbf{Ranking de equipos} En una competencia, compuesta por $T$ equipos, se enfrentan en cada uno de $P$ encuentros $2$ de los equipos. En cada partido, exactamente uno de los $2$ resulta ganador, por lo que no hay empates. El objetivo es determinar un orden entre los equipos, que refleje el desempeño de los mismos.

Para entender los distintos métodos posibles, veamos algunos ejemplos.

\begin{ejemplo}
	\normalfont
	Supongamos que en una competencia hay $6$ equipos, y que se han jugado $10$ partidos, en donde cada equipo anota goles y el ganador del partido resulta aquel que más goles convirtió, y no es posible el empate. Los resultados se pueden observar en la siguiente tabla:

	\begin{table}[h]
		\centering
		\begin{tabular}{|c|c|c|c|}
			\hline
			\textbf{Equipo local} & \textbf{Goles del local} & \textbf{Equipo visitante} & \textbf{Goles del visitante} \\ \hline
			1                     & 16                       & 4                         & 13                           \\ \hline
			2                     & 38                       & 5                         & 17                           \\ \hline
			2                     & 28                       & 6                         & 23                           \\ \hline
			3                     & 34                       & 1                         & 21                           \\ \hline
			3                     & 23                       & 4                         & 10                           \\ \hline
			4                     & 31                       & 1                         & 6                            \\ \hline
			5                     & 33                       & 6                         & 25                           \\ \hline
			5                     & 38                       & 4                         & 23                           \\ \hline
			6                     & 27                       & 2                         & 6                            \\ \hline
			6                     & 20                       & 5                         & 12                           \\ \hline
		\end{tabular}
	\end{table}
	
	
\end{ejemplo}
 

\newpage
\section{Resolución}


\newpage
\section{Experimentación}

\newpage
\section{Conclusiones}


\newpage

\end{document}